\documentclass[12pt]{amsart}
\usepackage[margin=1in,lmargin=0.5in,rmargin=1.5in]{geometry}
\usepackage{amsmath,latexsym,amssymb,amsthm,graphicx}
\usepackage{mathtools}
\usepackage{subfig}
\usepackage{xfrac}
\usepackage{fix-cm}
\usepackage{color}
\usepackage{verbatim}
\DeclareMathOperator\arctanh{arctanh}
\DeclareMathOperator{\arccosh}{arccosh}
\usepackage[bookmarks=true,%
    colorlinks=true,%
    linkcolor=blue,%
    citecolor=blue,%
    filecolor=blue,%
    menucolor=blue,%
    urlcolor=blue]{hyperref}
\usepackage{doi}
\usepackage{stmaryrd}
\usepackage{multirow}
\usepackage{float}
\usepackage{tabularx}
\usepackage[shortlabels]{enumitem}
\usepackage{pdfsync}
\usepackage{tikz}
\usepackage{tikz-cd}
\usepackage{float}
\usepackage[utf8]{inputenc}
\usepackage[textsize=small]{todonotes}
\newcommand{\DPTtodo}[2][]{\todo[color=green!40,#1]{#2}}
\newcommand{\DMGtodo}[2][]{\todo[color=red!40,#1]{#2}}
\usepackage{hyphenat}
\usepackage{booktabs}

\graphicspath{{mpdraws/}{draws/}}

\setenumerate[1]{label=(\arabic*)}

\numberwithin{equation}{section}
\numberwithin{figure}{section}

\newtheorem{theorem}{Theorem}[section]
\newtheorem{lemma}[theorem]{Lemma}
\newtheorem{claim}[theorem]{Claim}
\newtheorem{corollary}[theorem]{Corollary}
\newtheorem{proposition}[theorem]{Proposition}
\newtheorem{exercise}[theorem]{Exercise}
\newtheorem*{ex*}{Exercise}

\theoremstyle{definition}
\newtheorem{remark}[theorem]{Remark}
\newtheorem{remarks}[theorem]{Remarks}
\newtheorem{example}[theorem]{Example}
\newtheorem{definition}[theorem]{Definition}
\newtheorem{question}[theorem]{Question}
\newtheorem{conjecture}[theorem]{Conjecture}
\newtheorem{notation}[theorem]{Notation}
\newtheorem{fact}[theorem]{Fact}

\newcommand\N{\ensuremath{\mathbb{N}}}
\newcommand\Z{\ensuremath{\mathbb{Z}}}
\newcommand\Q{\ensuremath{\mathbb{Q}}}
\newcommand\R{\ensuremath{\mathbb{R}}}
\newcommand\C{\ensuremath{\mathbb{C}}}
\def\H{\ensuremath{\mathbb{H}}} % Redefining "long Hungarian umlaut"
\newcommand\D{\ensuremath{\mathbb{D}}}
\newcommand\F{\ensuremath{\mathcal{F}}}
\newcommand\G{\ensuremath{\mathcal{G}}}
\newcommand\J{\ensuremath{\mathcal{J}}}
\newcommand\U{\ensuremath{\mathcal{U}}}
\newcommand\V{\ensuremath{\mathcal{V}}}
\DeclareMathOperator{\Sp}{Sp}
\DeclareMathOperator{\Curr}{Curr}
\DeclareMathOperator{\Teich}{Teich}
\newcommand{\pa}[1]{\left(#1\right)}
\newcommand{\cb}[1]{\left\{#1\right\}}
\newcommand{\br}[1]{\left[#1\right]}
\newcommand{\rest}[1]{\left.#1\right|}
\newcommand{\card}[1]{\left| #1 \right|}
\newcommand{\ab}[1]{\langle #1 \rangle}
\renewcommand{\card}[1]{\tn{card}\pa{#1}}
\renewcommand{\Re}{\operatorname{Re}}
\renewcommand{\Im}{\operatorname{Im}}
\DeclareMathOperator{\ang}{ang}
\DeclareMathOperator{\Arg}{Arg}
\DeclareMathOperator{\Log}{Log}
\DeclareMathOperator{\len}{len}
\DeclareMathOperator{\tr}{tr}
\DeclareMathOperator{\diam}{diam}
\DeclareMathOperator{\EL}{EL}
\DeclareMathOperator{\Area}{Area}
\DeclareMathOperator{\interior}{int}
\DeclareMathOperator{\id}{id}
\newcommand{\del}{\partial}
\newcommand{\Res}[2]{\tn{Res}\br{#1, #2}}
\newcommand{\OneHalf}{{\textstyle \frac{1}{2}}}
\newcommand{\wt}[1]{\widetilde{#1}}

\renewcommand{\epsilon}{\varepsilon}

% Linear groups
\DeclareMathOperator{\ISO}{\mathit{ISO}}
\DeclareMathOperator{\SO}{\mathit{SO}}
\DeclareMathOperator{\GL}{\mathit{GL}}
\DeclareMathOperator{\SL}{\mathit{SL}}
\DeclareMathOperator{\PSL}{\mathit{PSL}}

% More delimiters
\DeclarePairedDelimiter\ceil{\lceil}{\rceil}
\DeclarePairedDelimiter\floor{\lfloor}{\rfloor}
\DeclarePairedDelimiter\abs{\lvert}{\rvert}
\DeclarePairedDelimiter\norm{\lVert}{\rVert}

% Specific notation for this paper
\newcommand{\Simples}{\mathcal{S}}
\newcommand{\Curves}{\mathcal{C}}
\newcommand{\Meas}{\mathcal{M}}
\newcommand{\PMeas}{\mathcal{M}_1}
\newcommand{\MSimples}{\Simples'}
\newcommand{\MCurves}{\Curves'}
\newcommand{\GC}{\mathcal{GC}}
\newcommand{\ML}{\mathcal{ML}}
\newcommand{\PML}{\mathcal{PML}}
\newcommand{\reducesto}{\mathrel{\searrow}}

\def\mathcenter#1{\vcenter{\hbox{$#1$}}}
\def\mfig#1{\mathcenter{\includegraphics{#1}}}
\newcommand{\svgfig}[1]{\Huge{\input{draws/#1.pdf_tex}}}
\newcommand*{\svgfignormal}[1]{\fontsize{12pt}{12pt}\selectfont\input{draws/#1.pdf_tex}}

\usetikzlibrary{arrows}
\tikzset{
    labl/.style={anchor=south, rotate=90, inner sep=.5mm}
}
\tikzstyle{every picture}=[> = to]
% Style for labels on arrows in commutative diagrams
\tikzset{cdlabel/.style={execute at begin node=$\scriptstyle,execute at end node=$}}
\tikzset{implication/.style={double equal sign distance, -implies}}
\tikzset{biimplication/.style={double equal sign distance, implies-implies}}

\begin{document}
\title{Characterizing intersection numbers}
\author[Martínez-Granado]{Dídac Martínez-Granado}
\address{Department of Mathematics\\
         Indiana University,
         Bloomington, Indiana 47405\\
         USA}
\email{dmartinezgr@gmail.com}

\author[Thurston]{Dylan~P.~Thurston}
\address{Department of Mathematics\\
         Indiana University,
         Bloomington, Indiana 47405\\
         USA}
\email{dpthurst@indiana.edu}
\date{Draft of \today}
\begin{abstract}
We give a characterization of functions defined on the space of closed curves that come from intersection numbers with a geodesic current. The key condition is that the function must be non-increasing under smoothing essential crossings.
\end{abstract}
\maketitle

\tableofcontents

\section{Introduction}

\begin{comment}
\textbf{Notation guidelines:}
\begin{enumerate}
\item  $S$ vs $\Sigma$: $S$ is a topological surface,
  $\Sigma$ has some extra structure like hyperbolic. Usually we want
  $S$ here
\item $a^{-1}$ vs $A$ for inverse element of group: Both are
  acceptable. Use $A$ if we're not specifically emphasizing the inverse
\item Use $g$ and other latin letters for group elements, and $\gamma$
  and other greek letters for geodesics of Poincare disk (or $\ora{g}$)
\item Figure colors for lifts of $a, b, ab, aB$, and so on should be uniformized.
\item Trees figures should be uniformized.
\end{enumerate}
\end{comment}
Geodesic currents on a hyperbolic surface are measures realizing a suitable closure of the space of weighted closed geodesics on the surface. Many metric
structures embed in the space of currents, such as
hyperbolic metrics \cite[Theorem~12]{Bonahon88:GeodesicCurrent},
 negatively curved Riemannian metrics
\cite[Th\'eor\`eme~1]{Otal90:SpectreMarqueNegative} or even non-positively curved Riemannian metrics
\cite[Theorem~A]{CFF92:RigidityNonPosCurvedRiem}.



Bonahon constructed
\cite[Proposition~4.5]{Bonahon86:EndsHyperbolicManifolds} a continuous bilinear function on the space of currents
extending the geometric intersection of closed multi-curves on the surface. 
Some notions of length of curves or \emph{curve functionals} $f$ can be realized as intersection numbers with a dual geodesic current $\mu_{f}$, this is,
\[
  i(\mu_{f},C)=f(C)
\]
for any closed curve $C$ in $S$. We will say that the current $\mu_{f}$ is \emph{dual} to the curve functional $f$, or that $f$ \emph{has a dual current} $\mu_{f}$. 
For example, the hyperbolic length of curves $\ell_{\Sigma}$ is dual to
a distinguished geodesic current~\cite[Theorem~10]{Bonahon88:GeodesicCurrent}), the \emph{Liouville current} $\mathcal{L}_{\Sigma}$; similarly for the negatively curved and non-positively curved Riemannian metrics mentioned above. Other examples do not come from metrics on the surface. For example, word-length with respect to simple generating sets have a dual current \cite[Theorem~1.2]{Erlandsson:WordLength}; or length functions arising naturally from certain representations of surface groups into higher rank Lie groups have corresponding dual currents (see~\cite{MZ19:PositivelyRatioed} and, more recently~\cite{BIPP21:PositiveCrossratios}).

\begin{comment}There are other natural notions of lengths associated to surface groups that are not dual to a geodesic current. For example, extremal length, or the hyperbolic length
of curves associated to a quasi-fuchsian (non-fuchsian) representation~\cite[Proposition~1.8]{BR24:Cubulations},\cite[Theorem A]{FF21:QF}.
\end{comment}
In this paper we give an axiomatic characterization of curve functionals that are dual to geodesic currents. 
Let $\mathbb{R}\mathcal{C}(S)$ denote the space of weighted multi-curves on $S$ (i.e. finite unions of non-negatively weighted non-necessarily-simple closed curves on $S$).

\begin{mainthm}
    Suppose that a curve functional $f \colon\mathbb{R}\mathcal{C}(S) \to \mathbb{R}$ satisfies the \emph{smoothing property}
     \[
      f\left(\mfig{curves-0}\right) \ge \max \left\{ f\left(\mfig{curves-1}\right), f\left(\mfig{curves-2}\right)\right\},
   \]
   for every essential crossing of a multi-curve on the surface and, moreover, $f$ is \emph{additive} on connected components, \emph{stable} and \emph{homogeneous} (see~\ref{subsec:mainhypotheses} for definitions).
Then, $f$ has a dual geodesic current, i.e., there exists $\mu_f$ so that
\[
i(\mu_f,C)=f(C)
\]
for every curve $C$.
\label{thm:intersect}
\end{mainthm}
See~Definition~\ref{def:essential-crossing} for essential crossing, which can be roughly defined as a crossing that cannot be removed by homotopy.

\begin{comment}
The injectivity of the map from the space of metrics to the space of geodesic currents is known as the marked length rigidity property.
Other examples of geometric structures embedding inside the space of geodesic currents are given by
Croke-Fathi-Feldman for non-positively curved Riemannian metrics
\cite[Theorem~A]{CFF92:RigidityNonPosCurvedRiem},
Hersonsky-Paulin \cite[Theorem~A]{HP97:RigidityNegCurvedCone} for
negatively curved metrics with conical singularities.
Duchin-Leininger-Rafi subsequently gave, as a subcase of the last one, a more explicit embedding for singular Euclidean structures
associated to quadratic differentials/half-translation structures
\cite[Theorem~4]{DLR10:DegenerationFlatMetrics}. Later, Bankovic-Leininger have embedded non-positively curved Euclidean cone metrics in the space of geodesic currents \cite[Main Theorem]{BL17:RigidityFlat}. The latest work in this direction is \cite{ELS21:Hyperbolic}. 

The way the embedding works in each of the above cases is by associating
a canonical geodesic current $\mathcal{L}_{\rho}$ to a given metric
structure $\rho$. This metric structure is usually defined up to a certain
equivalence---as happens, for example, for Teichm\"uller space (see
Theorem~\cite[Theorem~10]{Bonahon88:GeodesicCurrent}).  Then one shows
that this assignment is in fact an injection. This canonical geodesic
current $\mathcal{L}_{\rho}$ was coined \emph{Liouville current} by
Bonahon \cite[Section~2]{Bonahon88:GeodesicCurrent} for the case of
hyperbolic structures, but the term has also been used more generally
for the geodesic currents associated to other metric structures (see,
for instance, \cite[Section~3.1]{EU18:LengthCurrents}).
These Liouville currents are tied to the geometric structure as follows.
By Theorem~\cite[Proposition~4.4]{Bonahon86:EndsHyperbolicManifolds}, weighted multi-curves on a surface embed in the space of geodesic currents.
Also, in the same paper Bonahon
\cite[Proposition~4.5]{Bonahon86:EndsHyperbolicManifolds} that extends
the notion of geometric intersection
number as a continuous bilinear function the space of pairs of geodesic currents.
 This allowed him to extend hyperbolic length to geodesic currents by
following the principle of realizing it as an intersection number with
a distinguished geodesic current.
Similarly, other notions of length of curves $\ell_{\rho}$ associated to the above mentioned metric structures $\rho$ can be realized as intersection numbers with the corresponding Liouville current $\mathcal{L}_{\rho}$, this is
\[
  i(\mathcal{L}_{\rho},C)=\ell_{\rho}(C)
\]
for any closed multi-curve $C$ in $S$. We will say that the current $\mu_{\rho}$ is \emph{dual} to $\ell_{\rho}$.

There are many other lengths with an associated dual geodesic current.
For example, Erlandsson shows that word-length with respect to simple generating sets of the fundamental group can also be realized as an intersection number with such a Liouville current \cite[Theorem~1.2]{Erlandsson:WordLength}, which, in this case, turns out to be a weighted multi-curve. 
In the context of representations of surface groups into higher rank lie groups (specifically, certain Anosov representations), work of Martone-Zhang~\cite{MZ19:PositivelyRatioed} and, more recently, Burger-Iozzi-Parreau-Pozzetti~\cite{BIPP21:PositiveCrossratios} provides geodesic currents dual to length functions arising from logarithms of ratios of eigenvalues of the representation.

There are other natural notions of lengths associated to surface groups that are not dual to a geodesic current. For example, Brody-Reyes~\cite[Proposition~1.8]{BR24:Cubulations} show, mostly using work of Fricker-Furman~ \cite[Theorem A]{FF21:QF}, that for a quasi-fuchsian representation $\rho$, that is, a deformation of a fuchsian representation of $\pi_1(S)$ inside of $\PSL(2,\mathbb{C})$ whose limit set is still homeomorphic to a circle, the natural notion of stable length $\ell_{\rho}$ associated to this representation is dual to a geodesic current if and only if the representation is fuchsian (i.e., the deformation is trivial).

In fact, the usual first step in verifying marked length rigidity property for a family metric structures is to associate, to a metric structure $\rho$, a geodesic current $\mu_{\rho}$ dual to the length function $\ell_{\rho}$.
In this paper, we give an answer to the following question.

\begin{question}\label{q:intersection}
Let $\mathbb{R}\mathcal{C}(S)$ denote the space of weighted multi-curves on $S$.
What are sufficient and necessary conditions for $f \colon\mathbb{R}\mathcal{C}(S) \to \mathbb{R}$ to have a dual geodesic current?
\label{que:mainquestion}
\end{question}

\subsection{Main results}

\begin{table}
    \begin{tabular}{@{}ll@{}}
      \toprule
      Notation & Meaning \\
      \midrule
         $S$ &  topological surface \\
         $\Sigma$ & Riemannian surface \\
         $\Curves(S)$ & unoriented multi-curves \\
         $\Curves^+(S)$ & oriented multi-curves \\
         $\partial_{\infty} S$ & boundary at infinity of $\tilde{S}$ \\
         $G(S)$ & oriented geodesics on $\tilde{S}$ \\
         $\R\Curves(S)$ & weighted unoriented multi-curves \\
         $\R\Curves^+(S)$ & weighted oriented multi-curves \\ 
         $\GC(S)$ & unoriented geodesic currents \\
         $\GC^+(S)$ & oriented geodesic currents \\
         $\gamma$ & concrete multi-curve on $S$ \\
         $C$ & multi-curve on $S$\\
         $a, A \in \pi_1(S)$ & $A$ is the inverse of $a$, sometimes also $a^{-1}$\\
         $\ora{g}$ for $g \in \pi_1(S)$ & the hyperbolic axis of $g$
    \end{tabular}
  \caption{Notation for the objects related to surfaces, curves, and
    geodesic currents.}\label{tab:object-notation}
\end{table}

The main result of this paper can be succintly phrased as follows.

\begin{mainthm}
    Suppose that $f \colon\mathbb{R}\mathcal{C}(S) \to \mathbb{R}$ satisfies the \emph{smoothing property}
     \[
      f\left(\mfig{curves-0}\right) \ge \max \left\{ f\left(\mfig{curves-1}\right), f\left(\mfig{curves-2}\right)\right\},
   \]
   for every essential crossing of a multi-curve on the surface and, moreover, $f$ is \emph{additive} on connected components, \emph{stable} and \emph{homogeneous} (see~\ref{subsec:mainresult} for definitions).
Then, $f$ has a dual geodesic current, i.e., there exists $\mu_f$ so that
\[
i(\mu_f,C)=f(C)
\]
for every multi-curve $C$.
\label{thm:intersect}
\end{mainthm}
See~Definition~\ref{def:essential-crossing} for essential crossing, which can be roughly defined as a crossing that cannot be removed by homotopy.\end{comment}

This can be seen as a duality result giving a one-to-one correspondence between geodesic currents  $\GC(S)$ on surfaces $f \colon \mathcal{C}(S) \to \mathbb{R}$ satisfying extra conditions. See Section~\ref{subsec:newdef}. By \cite[Theorem~A]{MGT21:Smoothings} we obtain a continuous extension of arbitrary Riemannian length as an intersection number with a current. This also follows from the fact that $i(\mu_f, \cdot)$ extends continuously to geodesic currents by Bonahon's theorem of continuous bilinear extension of geometric intersection number (see \cite[Proposition~4.5]{Bonahon86:EndsHyperbolicManifolds}).
    A continuous extension of an arbitrary Riemannian length was previously given by \cite[Theorem~1.5]{EPS:CountingCurves}. However, as a consequence of Theorem~\ref{thm:intersect} we obtain a stronger result: a continuous extension as an intersection number. In fact, this also works more generally for length w.r.t any length metric structure on $S$ (for example, for any Finsler metric). As a further application of this fact, we show in Section~\ref{subsec:ineqintersection} that domination of intersection numbers with all closed curves implies domination of length spectra for any metric length on a closed surface.

    The smoothing condition as stated is easier to check for curve functionals arising from length or geodesic structures on the surface.
    For curve functionals arising from algebraic settings, such as the stable length of group actions on real trees (Section~\ref{sec:skora}), or cross-ratios (Section~\ref{sec:crossratios}), it is more convenient to decompose the smoothing condition into cases.

  Let $a,b \in \pi_1(S)$ and let
  $\ora{a}, \ora{b}$ be the oriented geodesics that are the corresponding hyperbolic axes (for some hyperbolic structure fixed once and for all).
 We say the pair $(\ora{a},\ora{b})$  \emph{R-cross} if $a^-, b^-, a^+,b^+$ appear in this
      ccw order on the boundary at infinity. See
      Figure~\ref{fig:cross}: we say the pair $(\ora{a},\ora{b})$ is \emph{R-parallel} if $b^-,b^+, a^+, a^-$ appear in this ccw order on the boundary at infinity. See Figure~\ref{fig:parallel}. One can similarly define a notion of $L$-crossing and $L$-parallel pair. Let a \emph{crossing} (resp. \emph{parallel}) pair mean $L$-crossing or $R$-crossing (resp. $L$-parallel or $R$-parallel). See Definition~\ref{def:crossing_parallel} for more details.
    
\begin{mainthmprime}
Let $a,b \in \pi_1(S)$.
Let $f \colon \pi_1(S) \to \mathbb{R}$ be a function satisfying the following conditions:
\begin{enumerate}
\item class invariant $f(a)=f(b)$ whenever $a$ is conjugate to $b$.
\item inverse invariant: $f(a)=f(A)$, for $A$ the inverse of $a$. 
\item stable: $f(a^n) = nf(a)$.
\item disconnected smoothing: For any pair of $R$-crossing elements $(a,b)$, 
 \begin{align}
    \label{eq:disconn_oriented} f(a) + f(b) &\geq f(ab)
  \end{align}


\item connected smoothing: For any pair of $R$-parallel elements $(a,b)$, 
 \begin{align}
    \label{eq:conn_oriented} \mbox{oriented} \quad f(ab) &\geq  f(a) + f(b)\\
    \label{eq:conn_unoriented} \mbox{unoriented} \quad f(ab) &\geq f(aB).
  \end{align}
\end{enumerate}
Then, there exists a geodesic current $\mu_f$ on $S$ so that
\[
i(\mu_f, [g]) = f(g)
\]
for every $g \in \Gamma$.
Moreover, any function of the type $i(\mu, \cdot)$, for a geodesic current $\mu$, satisfies the above assumptions.
\label{thm:intersections_group}
\end{mainthmprime}
See subsection~\ref{subsec:algebraic_vs_top} for a proof of the equivalence of the smoothing condition in Theorem~\ref{thm:intersect} to conditions (4) \emph{and} (5) in Theorem~\ref{thm:intersections_group}.
We emphasize that condition (4) alone is not enough to imply the smoothing condition: see Remark~\ref{rmk:reducible_action} and also Example~\ref{ex:linEL}.

One should compare this result to previous work of Feng
Luo~\cite{Luo98:IntersectionNumbers}, where he characterizes curve
functionals on simple closed curves which can be realized as
intersection numbers with measured laminations, i.e., geodesic currents $\alpha$ for which $i(\alpha,\alpha)=0$ (see \cite[Proposition~17]{Bonahon88:GeodesicCurrent}). Luo obtained a 
characterization of measured laminations in terms of simple curve
functionals~\cite[Theorem~1.1,Theorem~1.2]{Luo98:IntersectionNumbers}. In Section~\ref{sec:appsintersection}, we use Theorem~\ref{thm:intersect} and an
argument about fundamental domains of intersecting curves (for the
$\pi_1(\Sigma)$-action on $G(\Sigma)$) to give a new characterization
of measured laminations in terms of curve functionals~\ref{thm:char_laminations}. Feng Luo also
characterized which functions on simple closed curves are dual to
hyperbolic lengths~\cite{Luo98:GeodesicLength}. Our
Theorem~\ref{prop:char_liouville} can be seen as an extension of that
theorem for closed curves. We give a further characterization in Theorem~\ref{thm:char_liouville_strong}, in terms of a condition on the curve functional which is a priori easier to check.
 We outline the relation between our conditions and Feng Luo's conditions in the Appendix, Proposition~\ref{prop:feng}.
 
 In a previous paper, \cite[Proposition 4.11]{MGT21:Smoothings}, we have proven that the stable length associated to a quasi-fuchsian representation satisfies \emph{quasi-smoothing}, i.e., smoothing up to uniform additive error (see Definition \ref{def:properties} for details). Recently, Brody-Reyes~\cite[Proposition~1.8]{BR24:Cubulations}, following essentially on work of Fricker-Furman \cite[Theorem A]{FF21:QF}, show that this additive error in the quasi-smoothing cannot be removed unless the representation is fuchsian. 
In the case the representation is fuchsian, by \cite[Proposition 14]{Bonahon88:GeodesicCurrent}, it is dual to a geodesic current, and hence smoothing is satisfied.
Our result implies that, in fact, the smoothing property is a \emph{necessary} condition for \emph{any} length dual to a geodesic current. More importantly, adding a few more mild assumptions, it shows that it is also a \emph{sufficient} condition.

\subsubsection{Characterize measured laminations and hyperbolic metrics}

As a corollary of the Theorem and further work, we can further characterize which curve functionals correspond to measured laminations or to hyperbolic Liouville currents.

\begin{mainthm}
Let $f$ be a curve functional satisfying the hypotheses of Theorem~\ref{thm:intersect}. Then,
\begin{itemize}
\item $\mu \in \Curr(X)$ is a measured lamination if and only if $f$ satisfies the so-called \emph{disconnected max-smoothing property}:
\begin{equation*}
      f\left(\mfig{curves-10}\right) = \max \left\{ f\left(\mfig{curves-11}\right),f\left(\mfig{curves-12}\right) \right\}.
    \end{equation*}
    \item $\mu \in \Curr(X)$ is a hyperbolic Liouville current if and only if the functional $f$ satisfies the following condition:\[
        \lambda \left(\mfig{curves-10}\right) = \lambda \left(\mfig{curves-11}\right) + \lambda \left(\mfig{curves-12}\right),
  \]
  where $\lambda( \cdot ) \coloneqq 2\cosh(f(\cdot)/2)$.
\end{itemize}
    The left-hand-side of both equations depicts an essential crossing between any pair of curves $C_1$ and $C_2$, and the right-hand-side shows the two possible resolutions.
    \label{mainthm:apps}
\end{mainthm} 

 In addition to assuming the disconnected max-smoothing hypothesis, one can further assume the \emph{connected max-smoothing} hypothesis, obtained by replacing the essential crossing between two closed curves by a self-crossing in the max-smoothing equation above. This property is also satisfied by curve functionals dual to measured laminations and, together with additivity, homogeneity and stability, characterizes curve functionals dual to measured laminations~(see Theorem~\ref{thm:char_laminations}).

Furthermore, the second family of equations characterizing Liouville currents can be replaced by a family of inequalities (which are, a priori, easier to check) providing an alternative characterization of Liouville currents (see Section~\ref{subsec:liouville}).

\subsubsection{Characterizing small actions on real trees}

As an additional application, we reprove a result of Skora, giving a characterization of actions of surface groups on $\mathbb{R}$-trees that are dual to measured laminations.
 Let $\Gamma$ be a surface group, i.e. $\Gamma=\pi_1(\Sigma)$ for some closed hyperbolic surface $\Sigma$, and $T$ an $\mathbb{R}$-tree. Suppose $\Gamma$ acts on $T$ by isometries. It turns out that an isometry $g$ of $T$ either has no fixed points and an invariant axis in $T_g \subset T$ (hyperbolic), or it has fixed points (elliptic).
The action \emph{preserves axes intersection} if every pair of isometries of the surface $g, h \in \Gamma$ with transversely intersecting axes in $\wt{\Sigma}$, and acting hyperbolically on $T$, have intersecting axes in the tree.
An example of such an action is an $\mathbb{R}$-tree $T_{\lambda}$ \emph{dual to a measured lamination} $\lambda$, which, roughly speaking, is obtained by taking the Hausdorff quotient of $\wt{\Sigma}$ with the pseudo-metric $d_{\lambda}(x,y)$ obtained by considering the measure of all geodesic transverse to the hyperbolic segment between the points $x, y \in \wt{\Sigma}$.
We give a new proof of the following result, that was originally proven by Skora in~\cite{Sko90:GeometricAction}.

\begin{mainthm}
A minimal action of $\pi_1(S)$ on an $\mathbb{R}$-tree $T$ is dual to a measured lamination if and only if it is irreducible and preserves axes intersection.
\label{mainthm:skora}
\end{mainthm}





Culler-Morgan~\cite{CM87:GroupActions} give axioms for translation lengths on $\mathbb{R}$-trees which relate to (but are weaker than) our smoothing condition. In fact, in Proposition~\ref{prop:small} we give an alternate characterization of measured laminations in terms of actions whose translation lengths satisfy some additional axioms obtained by rephrasing the smoothing condition in the context of translation length of $\mathbb{R}$-trees.
Walter Parry~\cite{Par91:Translation} characterized functions on
conjugacy classes in $\Gamma$ that are translation lengths of $\Gamma$-trees in
terms of the Culler-Morgan axioms. One can see our result in terms of the
smoothing condition as further characterizing which functions on
conjugacy classes in $\Gamma$ are translation lengths of \emph{small} $\Gamma$-tree actions (i.e. actions with cyclic edge stabilizers), when
$\Gamma=\pi_1(S)$. 

We prove the following.

\begin{mainthm}
Let $\pi_1(S)$ act on a real tree $T$ minimally, and let $\ell \colon \pi_1(S) \to \mathbb{R}$ be its associated translation length on the tree.
The action $\pi_1(S) \acts T$ is irreducible and satisfies axis intersection if and only if $\ell$ satisfies the following conditions.
\begin{enumerate}
\item For any pair of $R$-crossing elements $(a,b)$, 
 \begin{align}
    \label{eq:small_disconn} \ell(a) + \ell(b) & = \max \{ \ell(ab), \ell(aB) \}.
  \end{align}
\item For any pair of $R$-parallel elements $(a,b)$, 
 \begin{align}
    \label{eq:small_conn}  \quad \ell(ab) &= \max \{ \ell(a) + \ell(b), \ell(aB)\}.
  \end{align}
\end{enumerate}
\label{thm:skora_maxsmoothing}
\end{mainthm}

It would be interesting to give similar conditions for 
small actions of a non-elementary hyperbolic group $\Gamma$. 

\subsubsection{Generalized cross-ratios}

Finally, we subsume previous work of Burger-Iozzi-Parreau-Pozzetti~\cite{BIPP21:PositiveCrossratios} on associating a geodesic current to every generalized cross-ratio. Every generalized cross-ratio induces a curve functional, its \emph{period}, giving a function
\[
\pi \colon \operatorname{CR}(S) \to \CurveFunct(S).
\]
On the other hand, \cite{BIPP21:PositiveCrossratios} associate, to every cross-ratio, a geodesic current,
\[
c \colon \operatorname{CR}(S) \to \GC(S).
\]
We show that period curve functionals (i.e., the image of the map $\pi$) satisfy the conditions of our main Theorem~\ref{thm:intersect} and thus induce geodesic currents.

\begin{mainthm}
There is a one-to-one correspondence between periods of generalized cross-ratios and geodesic currents.
\end{mainthm}

Unlike the work in~\cite{BIPP21:PositiveCrossratios}, our construction relies solely on the period function to generate the geodesic current. Although a given geodesic current is associated with two distinct cross-ratios, they share the same period. Our proof demonstrates that the mapping from cross-ratios to geodesic currents is mediated by the period map.







\subsection{Acknowledgments}
\label{sec:acknowledgments}

We would like to thank Giuseppe Martone, Beatrice Pozzetti, Eduardo Reyes, Noelle Sawyer, Juan Souto and Ivan Telpukhovskiy for useful conversations.

The first author was supported by the Hazel King Thompson fellowship of the Mathematics Department of the Indiana University Bloomington, for the initial part of this work, and by the Luxembourg National research Fund AFR/Bilateral-ReSurface 22/17145118.
The second author was supported by the National Science
Foundation under Grant Numbers DMS-1507244 and DMS-2110143.

\section{Further questions}
\label{sec:questions}

We finish by discussing some natural follow-up questions and conjectures.

\subsection{Open surfaces}
A natural question is whether our results extend to open surfaces, i.e.\ compact surfaces with boundary or finite-type surfaces with punctures.

For such surfaces, geodesic currents admit at least two natural definitions, depending on how one treats the ends. 
If the ends are realized as cusps, one obtains a space $\Curr_{\mathrm{cusp}}(S)$, defined analogously to Definition~\ref{def:currents} as the space of $\pi_1(S)$-invariant 
Radon measures on the space of geodesics in the universal cover (see, e.g.,~\cite[Definition~2.2]{Sas22:Currents}). In this setting, the space of geodesics includes geodesic arcs connecting cusps. 
If instead one restricts to geodesics projecting to the convex core,
one obtains a space denoted $\Curr_{\mathrm{open}}(S)$ (\cite[Section~2.6]{DLR10:DegenerationFlatMetrics}). See
Figure~\ref{fig:open_surfaces}.

\begin{figure}
\centering
\begin{subfigure}{0.48\textwidth}
\centering
\resizebox{60mm}{!}{
\svgfigbig{cusp}}
\caption{Cusped surface with an arc $A$ and a closed geodesic $C_n$
  wrapping around the arc. We have $C_n \to 2A$.}
\end{subfigure}
\hfill
\begin{subfigure}{0.48\textwidth}
\centering
\resizebox{60mm}{!}{
\svgfigbig{open}}
\caption{Open surface with two boundary parallel curves $B_1$ and $B_2$. The sequence of curves $C_n$ also sits in this surface, and we have $\frac{1}{n}C_n \to B_1 + B_2$.}
\end{subfigure}
\caption{Examples of cusped and open surfaces, and currents on them.}
\label{fig:open_surfaces}
\end{figure}

In $\Curr_{\mathrm{cusp}}(S)$, the intersection pairing fails to be
continuous; see~\cite[Remark~2.26]{MGT21:Smoothings} and
also~\cite{Sas22:Currents, T22:ThurstonCompactification}. It is thus
hard to get started working entirely in this setting.


The case of $\Curr_{\mathrm{open}}(S)$ is different. Even though intersection number is continuous in this case, it is not true that our defining properties of $\AConvex(S)$ guarantee a dual geodesic current.
The key issue is that the intersection number is degenerate: any
geodesic current $\mu$ supported on boundary curves will have zero
intersection number with all geodesic currents $\nu \in \Curr(S)$.

To better understand the relation between the cusped and open
situation, consider Figure~\ref{fig:open_surfaces} in a surface with
two ends for simplicity.
In the cusped surface, let $A$ be the simple arc going from cusp to cusp, which is in $\Curr_{\mathrm{cusp}}(S)$.
Let $C_n$ be a sequence of closed geodesics obtained by following $A$
from near one cusp to the other, wrapping around the second cusp $n$ times,
going back along $A$, and wrapping around the first cusp $n$ times
before closing up. In the weak$^*$-topology of
$\Curr_{\mathrm{cusp}}(S)$, we have $\lim_{n\to\infty}C_n = 2A$. (See, for example,
Sasaki~\cite[Section~6]{Sas22:Currents}.)

In the $\Curr_{\mathrm{open}}$ case, let $B_1$ and $B_2$ be the
boundary parallel curves homotopic to each respective boundary
component. The sequence of curves $C_n$ can also be considered in
$\Curr_{\mathrm{open}}(S)$. Without scaling, $C_n$ does not converge
to any current, but it does projectively. In fact, $\lim_{n \to
  \infty} \frac{1}{n} C_n = B_1 + B_2$.

Consider the geometric intersection number between arcs and curves in the cusped surface and curves in the open surface, i.e., $i \colon \Curves_{\mathrm{cusp}}(S) \times \Curves_{\mathrm{open}}(S)$.
We claim
\[
i(C_n, \tfrac{1}{n} C_n) \to i(2A, B_1 + B_2).
\]
Indeed, $i(C_n, C_n) = 4n$, and $i(2A, B_1+B_2)=4$, from which the result follows.
Following this example, we conjecture there is a well-defined and
continuous geometric intersection number between
$\Curr_{\mathrm{open}}(S)$ and $\Curr_{\mathrm{cusp}}(S)$.

\begin{conjecture}
There exist a continuous function 
\[
I \colon \Curr_{\mathrm{cusp}}(S) \times \Curr_{\mathrm{open}}(S) \to \mathbb{R}
\]
extending the natural geometric intersection pairing between $\Curves_{\mathrm{cusp}}(S)$ and $\Curves_{\mathrm{open}}(S)$.
\end{conjecture}

It is known that if one restricts the right-hand-side
factor to currents supported on a preimage of a compact subset on the
surface, the intersection number is continuous
(See \cite[Theorem~2.4]{BIPP21:Systoles}
and~\cite[Section~6]{Sas22:Currents}).

We can use this intersection number to propose a duality. Say that a
function $f \colon \Curves_{\mathrm{cusp}}(S) \to \mathbb{R}$
satisfies \emph{smoothing} if it satisfies smoothing between any
closed multi-curves, as well as smoothing between arcs and curves and
between arcs, in the sense that after resolving any essential crossing
(any intersection that cannot be homotoped away), the function does
not increase.

Similarly, say that a function
$f \colon \Curves_{\mathrm{open}}(S) \to \mathbb{R}$ satisfies
smoothing if it does not increase after resolving any essential
crossing.

For examples of such functionals, consider an embedded graph $\mathcal{G} \subset S'$.
Equipping $\mathcal{G}$ with a metric, we get a notion of length of
$\ell_\Sigma$ curves in $\Curves_{\mathrm{open}}(S)$ with respect to
this embedded metric graph. (See also
Section~\ref{subsec:embedded_graphs}.) When $\mathcal{G}$ is a rose-graph
(embedded and
with one vertex), Erlandsson~\cite{Erlandsson:WordLength} has
constructed examples of geodesic currents $\mu_{\mathcal{G}}$ in
$\Curr_{\mathrm{open}}(S)$ (in fact, half-integer weighted arcs and
multi-curves) so that
\[
\ell_{\Sigma}(C)=I(\mu_{\mathcal{G}}, C).
\]
This motivates the following question.

\begin{question}\label{quest:open-to-cusp}
Given a function $f \colon \Curves_{\mathrm{open}}(S) \to \mathbb{R}$
which is additive and satisfies smoothing, are there additional
assumption guaranteeing the existence of a unique geodesic current $\mu \in f \colon \Curr_{\mathrm{cusp}}(S)$ so that
\[
f(C)=i(C, \mu)
\]
for every $C \in \Curves_{\mathrm{open}}$?
\end{question}

Unfortunately, additivity and smoothing alone are not enough. For
example, take the function
$f \colon \Curves_{\mathrm{open}}(S) \to \mathbb{R}$ assigning
hyperbolic length to non-boundary parallel curves and $0$ to
boundary-parallel curves. This functional still satisfies smoothing
(since no other curve can cross a boundary-parallel curve essentially)
but does not have a dual geodesic current in
$\Curr_{\mathrm{cusp}}(S)$, since by continuity of $I$ it would have
to assign positive length to limits of curves wrapping around the
boundary. See~\cite[Remark~2.26]{MGT21:Smoothings} for details.

In the other direction, we have a parallel conjecture.

\begin{conjecture}\label{conj:cusp-to-open}
Given a function $f \colon \Curves_{\mathrm{cusp}}(S) \to \mathbb{R}$ which is additive and satisfies smoothing, there exists a unique geodesic current $\mu \in f \colon \Curr_{\mathrm{open}}(S)$ so that
\[
f(C)=i(C, \mu)
\]
for every $C \in \Curves_{\mathrm{cusp}}$(S).
\end{conjecture}

We have two motivating examples of functionals for
Conjecture~\ref{conj:cusp-to-open}.

The first one is related to the example on graphs above. Take again an
embedded metric graph $\mathcal{G} \subset S$, and define the
intersection $i(\mathcal{G},C)$ to be the minimal transversal
intersection between homotopy classes of embedded representatives of
$\mathcal{G}$ into $S$ and a homotopic representative of $C$, for
every closed curve in $\Curves_{\mathrm{cusp}}(S)$. This defines a
function $f \colon \Curves_{\mathrm{cusp}}(S) \to \mathbb{R}$. Then
our conjecture would assert the existence of a geodesic current
$\mu \in \Curr_{\mathrm{open}}(S)$. Presumably, as in Erlandsson's
result, $\mu$ in this case consists of half-integral combination of
curves.

As a second example, consider a flat metric on the completion of $S'$ so that the missing points become conical singular points. Now the length of an arc between cusps is the length of the arc from cone point to cone point, which is a finite number.
This gives a well-defined notion of length $\ell \colon \Curves_{\mathrm{cusp}}(S) \to \mathbb{R}$.
On the other hand,  $\Curr_{\mathrm{open}}(S)$ are geodesic currents
whose support is contained in the closure of the subspace of geodesics
that do not go through cone points. The natural Liouville metric for
singular flat metrics is an example of such a
current~\cite[Section~3.2]{BL17:RigidityFlat},
and seems to be dual
to~$\ell$.


\subsection{Generalizations to hyperbolic groups}
In recent work of Cantrell and Reyes and, independently, Kapovich and the first author, the following property has played a key role in extending stable lengths coming from group actions of hyperbolic groups $\Gamma$ acting on uniformly quasi-geodesic hyperbolic spaces $X$~\cite{KMG:BBP,CR25:Manhattan, CMGR26:GreenMetrics}.
Roughly, a Gromov hyperbolic space is a metric space whose triangles are thin. A group $\Gamma$ is said to be Gromov hyperbolic if its Cayley graph is.

A metric space $X$ is uniformly quasi-geodesic if there exists
constants $\lambda >0$ and $\epsilon \geq 0$, so that
every two points $x, y \in X$ are connected by a
$(\lambda,\epsilon)$-quasi-geodesic. We say that an action
$\Gamma \acts X$ of a non-elementary hyperbolic group~$\Gamma$ on a uniformly
quasi-geodesic hyperbolic metric space~$(X,d)$ has the \emph{bounded
  backtracking property} (see~\cite[Section~6.1]{CR25:Manhattan}) if given an orbit map
$\iota \colon \Cay(\Gamma) \to X$, there exists a constant $C>0$ so that
for every two points $x,y \in \Cay(\Gamma)$, and geodesic $[x,y]$,
there exists some quasi-geodesic $[\iota(x),\iota(y)]$ in $X$ so that
$\iota([x,y])$ and $[\iota(x),\iota(y)])$ are $C$-Hausdorff close in $d$.

There are examples of stable lengths satisfying bounded backtracking
which do not satisfy quasi-smoothing
(see~\cite[Section~9]{CMGR26:GreenMetrics}). On the other hand, one
can ask if all functions on $\pi_1(S)$ satisfying quasi-smoothing
arise as the stable length of some action of $\pi_1(S)$ satisfying
bounded backtracking. 

\begin{question}
Let $\Gamma$ be a surface group, and let $f \colon \Gamma \to \mathbb{R}$ be a symmetric and stable functional satisfying quasi-smoothing (and hence, extending continuously to geodesic currents).
Does there exists a uniformly quasi-geodesic Gromov hyperbolic space $X$ on which $\Gamma$ acts with the bounded backtracking property and so that $\ell_X = f$, where $\ell_X$ denotes the stable length of the action $\Gamma \acts X$?
\end{question}

Motivated by this, it would be interesting to generalize the quasi-smoothing and smoothing condition to the general setting of hyperbolic groups.

\begin{question}
What is a generalization of the properties of smoothing and
quasi-smoothing for functionals on non-elementary hyperbolic groups
that makes the answer to the
above question work in that setting?
\end{question}

Related to Theorems~\ref{thm:intersections_group}
and~\ref{mainthm:skora}, it would be interesting to give a list of
axioms, in the spirit of oriented smoothing, characterizing
small actions of a non-elementary hyperbolic group $\Gamma$.

\begin{question}
\label{qu:generalize_intersections}
Let $\Gamma$ be a non-elementary Gromov hyperbolic group action on an
$\mathbb{R}$-tree $T$. Give axioms  characterizing when the action is small, i.e., has
virtually cyclic edge stabilizers.
\end{question}

\subsection{Quasi-morphisms} 
\label{sub:quasimorphism} 

The hypotheses of \emph{oriented quasi-smoothing} are reminiscent of, but weaker than, the conditions for a map 
$f \colon \pi_1(S) \to \mathbb{R}$ to be a \emph{quasi-morphism}, that is, there exists $D > 0$ such that 
\[
|f(ab) - f(a) - f(b)| < D \quad \text{for all } a,b \in \pi_1(S).
\]
A quasi-morphism is called \emph{homogeneous} if $f(a^n) = n f(a)$ for every $a \in \pi_1(S)$ and $n \ge 1$.

Quasi-morphisms have played a prominent role in the study of stable commutator length and bounded cohomology of groups~\cite{Cal09:scl,Pic97:BoundedCohomology_currents}.

It would be interesting to understand whether one can obtain homogeneous quasi-morphisms from curve functionals satisfying oriented quasi-smoothing and stability.

\begin{question}
Let $f \colon \pi_1(S) \to \mathbb{R}$ satisfy oriented quasi-smoothing and stability.  
Are there additional conditions on $f$ ensuring that the function
\[
g(a) \coloneqq f(a) - f(a^{-1})
\]
is a homogeneous quasi-morphism?
\end{question} 

To get a non-trivial function~$g$, the functional~$f$ must be
asymmetric. See Appendix~\ref{sec:asymmetric} for examples.

% \subsection{Functionals valued on ordered abelian groups}

% We expect most of the results in this paper to extend, with only minor modifications, to the setting of $\Lambda$-valued functionals and $\Lambda$-valued geodesic currents, where $(\Lambda, <)$ denotes an ordered abelian group---that is, an abelian group equipped with a total order satisfying $a < b$ and $c \le d$ imply $a + c < b + d$.
% This generalization is natural in the context of the theory of $\Lambda$-trees~\cite{Sko90:GeometricAction} and has already proved fruitful in the study of compactifications of character varieties of surface-group representations into higher-rank Lie groups~\cite{BIPP24:PositiveCrossratios}.
% We plan to explore these extensions and their applications in future work.

\subsection{Green metrics}
\label{subsec:green}
Let $\Gamma=\pi_1(\Sigma) \subset \Isom^+(\mathbb{H})=\PSL_2(\mathbb{R})$ be a countable, discrete cocompact subgroup, for some hyperbolic metric $\Sigma$ on $S$.
Let $\mu$ be a finitely supported probability measure on $\Gamma$ whose support contains a symmetric generating set of $\Gamma$.
Let $w_n \coloneqq g_1 \cdots g_n$ be the random walk with respect to $\mu$, where each $(g_i)$ is identically and independently distributed with distribution $\mu$. Let $(\Gamma^{\mathbb{N}}, \mu^{\mathbb{N}})$ be the step space, and the map $\pi \colon \Gamma^{\mathbb{N}} \to \Gamma^{\mathbb{N}}$ as $\pi((g_n)_{n \in \mathbb{N}}) \coloneqq (w_n)_{n \in \mathbb{N}}$. The image of $\pi$ is equipped with the measure $(\mathbb{P}_{\mu})_*(\mu^{\mathbb{N}})$.
Define the metric \[ d_{\mu}(x,y) \coloneqq - \log \mathbb{P}_{\mu}(\exists n : w_n x =y).\]
In words, we are considering the probability that a random walk departing from $x$ will reach $y$ after some number of steps.
This metric has an associated stable length function $\ell_{\mu}(g) \coloneqq \lim_n d_{\mu}(x, g^n x)/n$, for any hyperbolic element $g \in \Gamma$, which can then be thought as a curve functional $\ell_{\mu} \colon \Curves^+(S) \to \mathbb{R}$. The following question emerged from conversations of the first author with Eduardo Reyes.

\begin{question}
Does $\ell_{\mu}$ satisfy the hypotheses of
Theorem~\ref{thm:intersect} for some admissible probability measure $\mu$ (and thus come from a geodesic current)?
\label{q:green-smoothing}
\end{question}

\begin{question}
Does $\ell_\mu$ further satisfy the
hyperbolic metric hypotheses in Theorem~\ref{thm:hyp-lam-char}(H) for some admissible
probability measure~$\mu$ (and thus come
from a hyperbolic metric)?
\label{q:Green}
\end{question}

If answers to either of these questions is negative for every $\mu$,
they would imply a negative answer to the following well-known question.

\begin{question}
Given $\Sigma$, is there a finitely-supported probability measure
$\mu$ on $\Gamma=\pi_1(\Sigma)$ so that the induced translation length
$\ell_{\mu}$ is dual to the Liouville current $\mathcal{L}_{\Sigma}$?
\label{q:singularity}
\end{question}

By work of Cantrell-Tanaka~\cite[Theorem~1.2]{CT25} and
Gouëzel-Mathéus-Maucourant \cite[Theorem~1.2]{GMM18:Entropy}, a
negative answer to this question would settle the so-called
\emph{singularity conjecture} for $\SL_2(\mathbb{R})$
(see~\cite[Conjecture, page~~259]{KLP11:Singularity}). In short, by
these results Question~\ref{q:singularity} is equivalent to asking if
there is a probability measure $\mu$ on $\Gamma$ so that the induced
`hitting measure' $\nu_{\mu}$ on $\partial \Gamma$ is absolutely
continuous with respect to the Lebesgue measure (coming from the
hyperbolic metric $\Sigma$) on $\partial \Gamma$.
Kosenko-Tiozzo~\cite{KT22:Inequality} have studied the metric $d_{\mu}$ in some cases of Fuchsian groups $\Gamma$ with centrally symmetric fundamental domains.
Even a positive answer to Question~\ref{q:Green} would be interesting, since then
Theorem~\ref{thm:intersect} would give a dual geodesic current for the
Green metric, and then one could study what types of geodesic currents
can arise.

\subsection{Smoothing vs multi-smoothing}
\label{subsec:multismoothing}
The hypotheses of smoothing in Equation~\eqref{eq:smoothing} of Theorem~\ref{thm:intersect} apply to one essential crossing at a time. There is a stronger condition than smoothing, where multiple smoothings at essential crossings are performed at once, and the function is required to be non-increasing under such operation.
We will refer to this condition as \emph{multi-smoothing}.
It is clear multi-smoothing implies smoothing.
It is easily checked that intersection numbers satisfy the multi-smoothing property (see, for example, \cite{AH21:EffectiveCounting3}).
What is not so clear is that functionals of $\AConvex(S)$ (that is, satisfying stability, homogeneity, smoothing and additivity) must satisfy multi-smoothing. This is, of course, a consequence Theorem~\ref{thm:intersect}.
Other curve functionals that are not in $\AConvex(S)$ also satisfy
multi-smoothing. For example, Arana-Herrera also shows
that extremal length satisfies
multi-smoothing \cite{AH21:EffectiveCounting3}.
This prompts the following question.

\begin{question}\label{quest:multi-smoothing}
If a curve functional satisfies homogeneity, stability, convex union,
and smoothing, does it necessarily satisfy multi-smoothing?
\end{question}

In this context, Arenas and Neumann-Coto
recently studied the related concept of \emph{taut smoothing}  \cite{ANC24:TautSmoothings}.

\subsection{Domination}
\label{subsec:domination}
In another direction from Section~\ref{subsec:multismoothing},
we say that a curve $C$ \emph{dominates} another curve $C'$ if $i(C,D) \geq
i(C',D)$ for every curve $D$. We say a functional $f$ satisfies
\emph{domination} if for
every pair of curves $(C,C')$ with $C$ dominating $C'$, we have
$f(C) \geq f(C')$.
It is easy to see that if $C'$ is obtained from $C$ by
multi-smoothing, then $C$ dominates $C'$. There are however other
pairs $(C,C')$ with $C$ dominating $C'$ but $C'$ not obtainable from
$C$ via multi-smoothing
(see~\cite[Example]{NC01:ShortestGeodesics}).
It is clear, by definition, that intersection numbers satisfy domination.
All in all, we have the following implications on a functional~$f$:
\[
  \text{Domination} \Rightarrow \text{Multi-smoothing} \Rightarrow \text{Smoothing}.
\]

From Theorem~\ref{thm:intersect} we obtain the following weaker
theorem which is nevertheless conceptually interesting.

\begin{theorem}
A curve functional satisfies domination and additivity if and only if it is dual to a geodesic current.
\end{theorem}

This can be interpreted as saying that linear maps that preserve the poset structure of the space of multi-curves correspond to intersection numbers.

For similar reasons as for the smoothing (or multi-smoothing)
condition, the additivity hypothesis is necessary, since
the square
root extremal length can be easily seen to satisfy domination,
homogeneity and stability, but not additivity, and it is not dual to
an intersection number.
This suggests the following strengthening of
Question~\ref{quest:multi-smoothing}.

\begin{question}
If a curve functional satisfies homogeneity, stability,
convex union, and smoothing, does it necessarily satisfy domination?
\end{question}

\subsection{Convexity on train-track charts and local compactness}
\label{sec:convexity_charts}

One of the reasons multi-smoothing is a useful property is that,
coupled with convex union, homogeneity, and stability, it
yields convex functions when restricted to train-track coordinates.
Precisely, let $f$ be a curve functional satisfying the aforementioned
properties, including multi-smoothing. Let $\tau$ be a maximal train-track (i.e., a
train-track all of whose complementary components are trigons), let
$V(\tau)$ be the set of admissible weights on~$\tau$, $\ML(\tau)$ be the subspace of
measured laminations carried by $\tau$, and
$\phi_{\tau} \colon V(\tau) \to \ML(\tau)$ be the induced
chart on measured laminations.
Since $f$ extends continuously to geodesic currents \cite[Theorem~A]{MGT21:Smoothings}, we can consider the composition $f_{\tau} \coloneqq f \circ \phi_{\tau} \colon V(\tau) \to \mathbb{R}$. Under these assumptions, the same proof as in~~\cite[Appendix~A]{Mir04:SimpleGeodesics}, \cite[Theorem~2]{Thu16:Research}, \cite[Theorem~4.7]{AH21:EffectiveCounting3}, shows $f_{\tau}$ is convex and homogeneous. 

These curve functionals form a subset of the functionals that extend
continuously to currents, by~\cite[Theorem~A]{MGT21:Smoothings}.
Moreover, the map $\Psi \colon \Curr(S) \to \mathbb{R}^{\Curves(S)}$
defined by $\Psi(\mu) = i(\mu, \cdot)$ is a homeomorphism onto its
image when $\mathbb{R}^{\Curves(S)}$ is equipped with the product
topology,
by~\cite[Theorem~11]{DLR10:DegenerationFlatMetrics}.
Thus, a restatement of Theorem~\ref{thm:intersect} is as follows.

\begin{theorem}
   The image of~$\Psi$ is homeomorphically $\AConvex(S) \subset \bbR^{\Curves(S)}$.
\end{theorem}

It is a consequence of this last result and the local compactness of
the space of geodesic currents (recalled as
Proposition~\ref{prop:local_compactness} below) that the space
$\AConvex(S)$ with the product topology is locally compact, which is
not obvious \textit{a priori} by looking directly at the functionals.

The sublevel set $\Lambda_f \coloneqq \{ \vec{x} \in V(\tau)
: f_{\tau}(\vec{x}) \leq 1 \}$ is a compact convex subset of
$\mathbb{R}^N$, for some $N>0$ depending on the combinatorics of~$\tau$.
For any sequence of geodesic currents $\mu_i$ limiting to~$\mu$,
let $f_i(\cdot) = i(\mu_i, \cdot)$ for $f(\cdot) = i(\mu, \cdot)$.
Then $\Lambda_{f_i} \to \Lambda_f$ in the Hausdorff topology. With
respect to this topology,
the subset of compact convex subsets of $\mathbb{R}^N$ is locally
compact \cite[Chapter~6]{Lay92:Convex}. We see this phenomenon as a
shadow of the local compactness of $\AConvex(S)$, which suggests that
the same local compactness for $\Convex(S)$ should hold.

\begin{conjecture}\label{conj:convex-compact}
The set $\Convex(S)$ with the product topology is locally compact.
\end{conjecture}

Some examples of interest for this conjecture would be extremal
lengths associated to conformal structures or to elastic graphs
\cite[Sections 4.8 and 4.9]{MGT21:Smoothings}.






\subsection{Analogy to Poincar\'e duality}
\label{subsec:poincare_duality}
Theorem~\ref{thm:intersect} gives a new definition of geodesic currents in terms of functionals on curves, by giving a bijection between geodesic currents and a subset of the space of curve functionals.
As mentioned in the introduction, one can view this theorem in analogy to Poincaré duality for closed surfaces.  Given a 1-cohomology class $\alpha^*$ of an $n$-dimensional manifold (represented by homomorphisms from $\pi_1(S)$ to $\mathbb{R}$), Poincaré duality associates an $(n-1)$-dimensional homology class $\alpha$ (represented by real linear combinations of $(n-1)$-cycles). Since $n=2$, these cycles can be expressed as linear combinations of simple closed curves.  

In our setting, the role of a 1-cohomology class is played by a curve
functional that satisfies smoothing, additivity, and
stability. Although it is not a homomorphism from
$\pi_1(S)$ to anything, it preserves a certain combinatorial structure
(see Section~\ref{subsec:domination}). Likewise, the role of the homology
class is taken by a geodesic current.

Our theorem can be thought of in analogy to $\alpha \cap \beta =
  \alpha^*(\beta)$, where $\cap$ denotes the cap product in homology (algebraic intersection), and thus relating cap product to the homology-cohomology pairing.

  This aligns with our theorem, which establishes a duality via the \emph{geometric} intersection number $i(\mu_f, C) = f(C)$, suggesting that the geometric intersection number plays the role of the cap product. 

  
In Poincaré duality, one also has the standard relation $\alpha^* \cup \beta^* = \alpha \cap \beta$, where $\cup$ denotes the cup product in cohomology.  
A natural question that arises then is the following.

\begin{question} What is the appropriate analog of the cup product?
  More precisely, can one define the intersection $i(\cdot,\cdot)$ of
  two geodesic
  currents directly via their dual curve functionals, and prove its
  continuity?
\end{question}

\subsection{Asymmetric analogues}
\label{sec:asymmetric-analogues}

In this paper, we characterized additive, unoriented curve functionals that
satisfy smoothing as exactly those coming from geodesic currents.
We could instead look at functionals on \emph{oriented} curves.

\begin{definition}
  Let $\AConvex^+(S)$ be the space of functionals on oriented curves
  $\Curves^+(S)$ that satisfy additivity, oriented smoothing
  (connected and disconnected), and stability, endowed with the
  product topology. Similarly, $\AConvex^+_{\ge 0}(S)$ is the subspace
  of positive functionals.
\end{definition}

Appendix~\ref{sec:asymmetric} can be summarized as saying that
$\AConvex^+(S)$ is not very related to the space $\Curr^+(S)$ of
oriented geodesic currents: the natural map $\mu \mapsto
\vec\imath(\mu,\cdot)$ is neither injective nor surjective, in what
appear to be large ways.

Nevertheless the space $\AConvex^+_{\ge 0}(S)$ may be of independent
interest. There are many examples of such functionals. For instance,
you can take an asymmetric metric on $S$ and obtain a functional
on curves as in Section~\ref{subsec:ineqintersection}. As in the
symmetric case, if there are rectifiable curves in each homotopy class
the result is automatically positive and satisfies oriented smoothing.

For a more concrete example, you can take an asymmetric generating set
for $\pi_1(S)$ as a monoid whose (directed) graph is simple in~$S$, as
in the example in Appendix~\ref{sec:asymmetric}.

\begin{question}\label{quest:AConvex+-compact}
  Is the space $\AConvex^+_{\ge 0}(S)$ a reasonable space? For
  instance, is it locally compact?
\end{question}

\begin{question}
  Does every element of $\AConvex^+_{\ge 0}(S)$ come from an
  asymmetric pseudo-metric in the style of
  Corollary~\ref{maincor:length_dual}?
\end{question}

In this oriented context, the restriction that the functional is
positive is non-trivial (unlike for unoriented functionals): for any
functional $f$ and element $\omega \in H^1(S; \bbZ)$, for a sufficiently
large constant $\lambda>0$ the combination $f + \lambda\omega$ will take negative values.

We could also relax the condition of smoothing to quasi-smoothing, and
ask similar questions. If we drop positivity and also consider
quasi-smoothings, then the space of functionals includes the much
larger space of quasi-morphisms (Section~\ref{sub:quasimorphism}), which an infinite dimensional Banach space and hence not locally compact (see, e.g.,~\cite{CRH26:Quasimorphisms}).

Recall that $\QConvex_K^+(S) \subset \QConvex^+(S)$ denotes the subspace of curve functionals on oriented multi-curves satisfying oriented quasi-smoothing with constant~$K$, homogeneity, stability and convex union.
$\QConvex_K^+(S)$ contains, for example, the space of homogeneous $K$-quasi-morphisms. The space of $K$-quasi-morphisms is easily seen to be locally compact by an inequality
\[
f(C) \leq K' \ell_F(C)
\]
where $\ell_F$ is the word-length induced by any finite generating set, and $K'$ is a constant depending only on a finite subset of curves.
The hypotheses of oriented $K$-quasi-smoothing are much weaker than those of $K$-quasi-morphisms.

\begin{question}\label{quest:QConvex+-compact}
Is $\QConvex_K^+(S)$ locally compact?
\end{question}

More generally, one could ask about local compactness for various weakenings of the the hypotheses defining the subset $\AConvex(S)$; Conjecture~\ref{conj:convex-compact} and Questions~\ref{quest:AConvex+-compact} and \ref{quest:QConvex+-compact} are some examples.

%%% Local Variables:
%%% mode: latex
%%% TeX-master: "Intersections"
%%% End:



%%% Local Variables:
%%% mode: latex
%%% TeX-master: "Intersections"
%%% End:

\section{Background}

\DMGtodo{Should talk a bit about Hamenstaedt and Ledrappier work, copy some material of background of curves From curves to currents here, talk about definition of geodesic current in terms of cross-sections, etc.}

\begin{comment}
\DMGtodo[inline]{The following section contains material that could be relevant to include in intro or background section if we end up using anything for applications.}
\section{A bunch of random background and previous ideas}

\textbf{Cocycles}

A class of H\"older continuous function $\zeta\colon T^1 S \times \mathbb{R} \to \mathbb{R}$ satisfying $\zeta(v,s+t)=\zeta(\phi_s(v),t) + \zeta(v,s)$
for all $v \in T^1 S$ and all $s,t \in \mathbb{R}$.
Let $\mathcal{F}\colon T^1 S \to T^1 S$ be the flip $v \mapsto -v$. The flip of a cocycle $\zeta$ is the cocycle defined by $\mathcal{F}\zeta(v,t)=\zeta(\mathcal{F}\phi^t(v),t)$. Two cocycles $\zeta, \xi$ are cohomologous if $\zeta(v,t)-\xi(v,t)=\beta(\phi^t(v)) - \beta(v)$ for all $v \in T^1 S$ and all $t \in \mathbb{R}$.
$\zeta,\xi$ are cohomologous iff $\mathcal{F}\zeta,\mathcal{F}\xi$ are cohomologous. Thus, the flip map induces a map in cohomology classes.
We say $\zeta$ is quasi-invariant under the flip if $\zeta$ and $\mathcal{F}\zeta$ are cohomologous.

\textbf{Cross-ratio}

A H\"older continuous, $\pi_1(S)$-invariant function on $\partial \wt{S}^{[4]}$ (space of quadruples of pairwise distinct points of $\partial \wt{S}$) with properties

\begin{itemize}
\item $Cr$ is invariant under the diagonal action of $\pi_1(S)$.
\item $Cr(a,a',b,b')=Cr(a',a,b,b')^{-1}$
\item $Cr(a,a',b,b')=Cr(b,b',a,a')$
\item $Cr(a,a',b,b')Cr(a',a'',b,b')=Cr(a,a'',b,b')$
\end{itemize}



\textbf{Periods}

\textbf{Periods for cross-ratio}
Given a homotopy class of a closed curve $c$ in $S$, we can consider a representative $\gamma$.
Let $\gamma$ be the Deck transformation associated to $\gamma$, and let $\gamma^-$ and $\gamma^+$ be its repelling and attracting points respectively.
The period $l(\gamma)=\log Cr(\gamma^+,\gamma^-,\gamma \cdot a, a)$ for any $a \notin \{\gamma^-,\gamma^+\}$.
This definition does not depend on the conjugacy class of the element $\gamma \in \pi_1(S)$, nor on the choice of point $a$. \\

\textbf{Periods for geodesic currents}

Let $c$ be a closed curve in $S$, lift it to the universal cover to some $\gamma$, choosing a hyperbolic metric. Let $G_{\gamma}$ be the set of geodesics that intersect $\gamma$ and have one endpoint in the segment $(a,\gamma a)$. $\nu(G_{\gamma})$ is the same as the measure of the set of geodesics that intersect the segment $(q,\gamma q)$ tranversally, for $q$ in the axis.




\textbf{Signed geodesic current}

A geodesic current for $S$ is a $\pi_1(S)$-invariant locally finite Borel measure on $\partial \wt{S} \times \partial \wt{S} - \Delta$ which in addition is invariant under the natural involution exchanging the two factors. Geodesic currents naturally correspond to finite $\phi^t$ invariant transverse Borel measures on $T^1M$ invariant under the flip map.
A signed geodesic current is a difference of two geodesic currents. \\

Otal proves that given a family of periods there is at most one cross ratio with those periods (essentially, density of endpoints of hyperbolic elements).
What are the properties a period must satisfy?

\textbf{From flip invariant Holder class to cross-ratio}

Let $\alpha$ be a cocycle representing the class. Let $x \in \wt{S}$ be any point and let $(,)_{\alpha}$ be the symmetric product on distinct points in $T^1_x \wt{S}$. For $a \neq b \in \partial \wt{S}$, write $(a,b)=(v,w)_{\alpha}$ where $v,w \in T^1 \wt{S}$ are such that $\pi(v)=a$ and $\pi(w)=b$. Then $(,)$ is a symmetric product on pairs of distinct points in $\partial \wt{S}$ depending on the choice of $x$ and $\alpha$. For a quadruple $(a,a',b,b')$ of distinct points in $\partial \wt{S}$ define $[a,a',b,b']_{\alpha}=(a,b)+(a',b')-(a,b')-(a',b')$. Then one checks that $[]_{\alpha}$ does not depend on the choice of $x$ or the representative $\alpha$. And $Cr=e^{[]_{\alpha}}$ is the cross-ratio associated to the class of $\alpha$.

\textbf{From cross-ratio to Holder cocycle}

One defines a cocycle for the frame flow in the frame bundle, then integrates over fibers to obtain a Holder cocycle.\\

Finally one checks that one is inverse of the other (if one closes the cycle obtains a cohomologous cocycle). \\

\textbf{From geodesic current to cross-ratio}

Let $\mu$ be a geodesic current. Let $(a,b,c,d)$ an ordered tuple of pairwise distinct points. Define $[a,b,c,d]=\mu([a,b] \times [c,d])$. \\
(To check Holder property need to use $\mu$ is Gibbs current and thus comes from a Holder function...?) \\
To check positivity one uses:

\begin{lemma}
Let $Cr$ be a cross-ratio on $\partial \wt{S}$. Then $Cr(a,b,c,d)+Cr(b,c,d,a) \to 1$ as $a \to b$, locally uniformly in $(c,d)$ if and only if $Cr$ is positive.
\end{lemma}

\textbf{From cross-ratio to intersection geodesic current}

For a generalized cross-ratio there is a unique signed current $\mu$. Define $\mu([a,b]\times[c,d])=\log Cr(a,b,c,d)$. Additivity of cross ratio gives additivity of measure. If the cross ratio is positive then it gives a geodesic current. \\

\textbf{From cocycle to Gibbs current}

Let $\mathcal{M}$ be the space of all $\phi^t$-invariant probability measures on $T^1 S$.

The entropy of a Borel set $A$ w.r.to to probability measure $\nu$ on $T^1S$ is $h_{\nu}(A):=-\nu(A)\log \nu(A)$. \\
The entropy of $\nu$ is the entropy of the whole space $T^1 S$. Let $\mathcal{S}$ be the $\sigma$-algebra of all measurable sets. The entropy of $\nu$ is

$h_{\nu} = \sup \{ \sum_{A\mathcal{F}} h_{\nu}(A) \mathcal{F} \subset \mathcal{S} | \mathcal{F} \mbox{ finite } , X=\sqcup_{A \in \mathcal{S}} A\}$. \\

For every H\"older continuous function $f\colon T^1 S \to \mathbb{R}$ the pressure $p(f):=\sup \{ h_{\nu} - \int f d\nu | \nu \in \mathcal{M} \}$.

The Gibbs equilibrium state of $f$ is the (unique) probability measure $\eta_f \in \mathcal{M}$ which realizes this supremum.

Any H\"older function $f\colon T^1 S \to \mathbb{R}$ gives rise to a H\"older cocycle by 

$$\xi_f(v,t):=\int_0^{t} f (\phi^s(v)ds.$$

We call $f$ quasi-invariant under the flip if for any periodic $v \in T^1 S$ of period $\tau$, $\xi_f(v,\tau)=\xi_f(-v,\tau)$. \\

For a H\"older cocycle, we can find some H\"older function $f\colon T^1 S \to \mathbb{R}$. Two H\"older $f,g$ functions give the same cocycle iff $\int fd\mu = \int g d\mu$ for all $\mu \in \mathcal{M}$.\\

Otal proves that cross-ratio is determined given by periods. \\

\subsection{Reconstructing currents from periods}

Otal's argument uses density of geodesic in tangent bundle to construct current in some rectangles. Idea: extend to all sigma-algebra using Caratheodory's theorem. Problems:

\begin{itemize}
\item Is it enough to reconstruct the current given only on this family of rectangles?
\item Continuity at endpoints.
\end{itemize}

\subsection{Reconstructing current}

Another approach is to try to define a positive functional on the space of compact continuous functions on the space of geodesic currents. \\
It also has to have the $\pi_1(S)$-invariant property.\\
The difficult point will be continuity of the functional. \\


Feng Luo characterizes functions that are intersection numbers on simple curves (ref).
\end{comment}

\subsection{Action of $\pi_1(S,*)$ on lifts}
\label{subsec:actionsfundamental}

There are two actions of $\pi_1(S,*)$ on lifts of $*$ to $\tilde{S}=\{ \mbox{paths from } * \}/\sim$:

\begin{enumerate}
    \item Left action: precomposition with $\gamma$ via Deck transformations, isometries of $\mathbb{H}^2$. The isometry extends to a homeomorphism of the boundary at infinity $S^1_{\infty}$. This action is minimal, in the sense that any orbit $\pi_1(S,*).p$, for $p \in S^1_{\infty}$, is a dense subset of $S^1_{\infty}$.
    \item Right action: post-composition of paths. For example, $a \in \pi_1(S,*)$ acts on $x \in \pi_1(S,*)$, as $ax$, where this is the path first running through $x$ and then through $a$. 
\end{enumerate}

%\begin{definition}
%  A homogeneous function on geodesic currents is \emph{convex} if its
%  restriction to curves is convex by Definition~\ref{def:convex}.
%\end{definition}

%\begin{definition}
%  \emph{Pants coordinates} on $\mathcal{ML}(S)$ are Dehn-Thurston
%  coordinates with respect to a marked pair-of-pants decomposition,
%  consisting of length and twist coordinates for each pants curve. The
%  twist parameters require more specification, since there are several
%  variations in the literature; see Section XXX for our version.
%\end{definition}

%\begin{theorem}
%  Let $f \colon \mathcal{GC}(S) \to \R_{\ge 0}$ be a convex
%  functional. Then the restriction of $f$ to measured laminations is
%  convex with respect to pants coordinates for any pair-of-pants
%  decomposition.
%\end{theorem}

%\subsection{Outline}
%Extending functions to ML: \\
%- Convexity argument, natural coordinates? \\
%- Examples:
%\begin{enumerate}
%\item Hyperbolic length (Thurston)
%\item Intersection number (Thurston)
%\item Extremal length (Kerckhoff)
%\end{enumerate}

%Extending functions to GC:
%\begin{enumerate}
%\item Hyperbolic length (Bonahon)
%\item Intersection number (Bonahon)
%\item Stable length (Bonahon, Erlandsson-Parlier-Souto)
%\item Generalized translation lengths (Martone-Zhang)
%\end{enumerate}

%In this paper we provide an argument that generalizes the previous extensions to a broad family of functions.

%Applications:

%\begin{enumerate}
%\item Counting: Rafi-Souto, Erlandsson-Parlier-Souto)
%\item Extending extremal length to geodesic currents, relevant in study of %rational maps.
%\item Others? (conversations with Wright).
%\end{enumerate}

%Inequalities:
%\[C_1 l(c) \leq f(c) \leq C_2 l(c).\]

%For general $f$, the right inequality follows from argument decomposing curve into generators (in the additive case).
%Question: can we also obtain it in the non-additive case, for general $f$? \\ Explain the case $f$ is extremal length. Consider the case of the torus. \\
%Right inequality follows from collar lemma, and other by considering embedded annulus, and looking at the core curve. \\
%These inequalities are interesting because: (1) they are satisfied by all Liouville currents (see paper by Burger et al on structure theorem for geodesic currents), and (2) the left inequality is usually a part of the definition of Anosov representation, where $f$ plays the role of the generalized translation length (see Martone-Zhang).

\subsection{Background on measure theory}

Let $X$ be a set.

\begin{definition}[Semi-ring, ring, $\sigma$-ring]
A family $\mathcal{S}$ of subsets of a set $X$ is called a \emph{semiring} if it contains the empty set, $A \cap B \in \mathcal{S}$ for all $A,B \in \mathcal{S}$ and, for every pair of sets $A,B \in \mathcal{S}$ with $B A \subset B$, the set $B \backslash A$ is a disjoint union of finitely many sets in $\mathcal{S}$.
A family $\mathcal{R}$ of subsets of a set $X$ is called a \emph{ring} if it contains the empty set and the sets $A \cap B$, $A \cup B$ and $A \backslash B$ belong to $\mathcal{R}$ for all $A,B \in \mathcal{R}$. 
The \emph{$\sigma$-ring} of $\mathcal{R}$ consists of the set of countable unions of elements of $\mathcal{R}$.
\end{definition}

\begin{definition}[semialgebra, algebra, $\sigma$-algebra]
If in the above definitions one assumes $X$ is also an element of the family, one gets, respectively, the definitions of \emph{semialgebra},\emph{algebra} and $\sigma$-algebra.
\end{definition}

\begin{lemma}
For any semiring $\mathcal{S}$, the collections of all finite unions of sets in $\mathcal{S}$ forms a ring $\mathcal{R}_{\mathcal{S}}$, the \emph{ring generated by $\mathcal{S}$}. Every set in $\mathcal{R}$ is a finite union of pairwise disjoint sets in $\mathcal{S}$.
\end{lemma}
\begin{proof}
See \cite[1.2.14~Lemma]{Bo07:MeasureTheory}.
\end{proof}

\begin{definition}
A real-valued set function $\mu$ defined on a class of sets $\mathcal{A}$ is called \emph{finitely additive} if 

\[
\mu\left(\bigcup_{i=1}^n A_i \right) = \sum_{i=1}^n \mu(A_i)
\]
for all $n$ and all disjoint sets $A_1,\cdots,A_n \in \mathcal{A}$ such that $\bigcup_{i=1}^n A_i \in \mathcal{A}$.
$\mu$ is called \emph{countably subadditive} if 

\[
\mu\left(\bigcup_{i=1}^{\infty} A_i \right) = \sum_{i=1}^{\infty} \mu(A_i)
\]
for all pairwise disjoint sets $A_i \mathcal{A}$ such that $\bigcup_{i=1}^{\infty} A_i \in \mathcal{A}$.
We say $\mu$ is called \emph{countably additive} if the above inequality is replaced by an equality.
\end{definition}

\begin{proposition}
A finitely additive function $\mu$ on a semiring $\mathcal{S}$ extends uniquely to a finitely additive function on $\mathcal{R}_{S}$
\label{prop:addsemitoring}
\end{proposition}
\begin{proof}
See \cite[1.3.10~Proposition]{Bo07:MeasureTheory}.
\end{proof}


\begin{definition}
  Fix a ring $\mathcal{R}$. Let $\mathcal{F}$ be a family of nested sequences
  of elements of $\mathcal{R}$. Then $\mathcal{F}$ is \emph{final} if,
  for any sequence $(R_i)_{i=0}^{\infty}$ of nested elements of
  $\mathcal{R}$ with $\bigcap_i R_i = \emptyset$, there is a sequence
  $(S_j) \in \mathcal{F}$ so that $R_i$ is eventually contained in
  $S_j$, in the sense that for each $j$ there is an $N$ so that for
  all $i \ge N$ we have $R_i \subset S_j$.
\end{definition}

To show countable additivity, suffices to show that there is a final
family on which measures go to zero.


\begin{proposition}
Let $\mu$ be a finitely additive function defined on a ring $\mathcal{R}$. Let $\mathcal{F}$ be a final family of $\mathcal{R}$ so that for any $(S_j)_{j=1}^{\infty} \in \mathcal{F}$, $\lim_j \mu(S_j) \to 0$. Then $\mu$ is countably subadditive.
\label{prop:finalfamcount}
\end{proposition}
\begin{proof}
Let $A_n$ be a sequence of pair-wise disjoint sets in $\mathcal{R}$ whose union $A$ is also in $\mathcal{R}$.
Then consider $B_n =  A - \cup_{i=1}^n A_i$ . The sequence $B_n$ is a nested decreasing sequence of elements in $\mathcal{R}$ with empty intersection. Since $\mathcal{F}$ is a final family, for every $j$ we can find $N$ so that for every $i \geq N$, we have $B_i \subset S_j$.
On the one hand, note that $\mu(B_i) \to 0$, since $\mu(S_j) \to 0$ by definition of final family and $\mu$ being monotone.
On the other hand, \[
\mu(A - \cup_{i=1}^n A_i )= \mu(A) - \mu( \cup_{i=1}^n A_i )
\]
by finite additivity of $\mu$, which proves what we wanted.
\end{proof}


\begin{definition}\DMGtodo{Is this same as outer regular in Radon definition? What about inner regular?}
A set function $\mu$ defined on the class of all subsets of a set $X$ and taking values in $[0,\infty]$ is called an \emph{outer measure} if

\begin{enumerate}
    \item $\mu(\emptyset)=0$
    \item $\mu(A) \leq \mu(B)$ whenever $A \subset B$, i.e., is monotone
    \item $\mu$ is countably subadditive. 
\end{enumerate}
A set $A \subset X$ is called $\mu$-measurable if, for every $E \subset X$, one has the equality

\[
\mu(E \cap A) + \mu(E \backslash A) = \mu(E)
\]
The class of all $\mu$-measurable sets is denoted by $\mathcal{M}_{\mu}$.
An outer measure is called \emph{regular} if, for every set $A \subset X$, there exist a $\mu$-measurable set $B$ so that $A \subset B$ and $\mu(A)=\mu(B)$.
\end{definition}

\begin{theorem}[Caratheodory extension theorem]
Given a ring $\mathcal{R}$ and a finitely additive and countably additive function $\mu$ on it.
Let 
\[
\mu^*(A) \coloneqq \inf \{ \sum_{i=1}^n \mu(A_n) | A_n \in \mathcal{R}, A \subset \cup_{i=1}^{\infty} X_i\},
\]
where in the case of absence of such sets $A_n$ we set $\mu^*(A) \coloneqq \infty$. Then

\begin{enumerate}
    \item  $\mu^*$ is an outer regular measure
    \item All the sets in $\mathcal{R}$ are $\mu$-measurable
    \item $\mu^*$ coincides with $\mu$ on $\mathcal{R}$
    \item $\mu$ extends to the countably additive extension to the $\sigma$-ring generated by $\mathcal{R}$
\end{enumerate}
If one replaces ring by algebra above, the extension in (4) is unique.
\label{thm:caratheodory}
\end{theorem}
\begin{proof}
See \cite[1.11.8~Theorem]{Bo07:MeasureTheory}.
\end{proof}

Let $X$ be a topological space.
Borel sigma algebra.
Define outer regular, inner regular, locally finite measures.
\section{Discussion on previous work}

By theorem \ref{thm:intersect} we get an analogue $\mu_g$ of the Liouville current. This current was constructed by Otal in \cite{Otal90:SpectreMarqueNegative} for all negatively curved Riemannian metrics. Then it was generalized in two different paths:
first by Croke-Fathi-Feldman \cite{CFF92:RigidityNonPosCurvedRiem} for non–positively curved Riemannian 
metrics and second by Hersonsky-Paulin \cite{HP97:RigidityNegCurvedCone}  for negatively  curved cone  metrics. The general case for arbitrary flat metrics appears to be new. \\
The existence of the analog of the Liouville current provides with the first step towards attempting the rigidity problem: is a family of metrics determined by its marked length spectrum? \\
By Otal \cite{Otal90:SpectreMarqueNegative} or our
Theorem~\ref{thm:intersect}, a geodesic current is determined by
intersection numbers with closed curves.
To determine spectral rigidity one has to show that two metrics inducing the same geodesic current are equivalent (up to maps isotopic to the identity). To do that one usually then studies the support of these geodesic currents. \\

Comment: current we associate to an arbitrary Riemannian metrics is
not the long-term behavior of a random geodesic. For instance, if $g$
has positive curvature, may have an open set of Riemannian geodesics
on~$\Sigma$ that stay inside a contractible set and so are not
quasi-geodesic on~$\wt \Sigma$.



\section{From additivity to geodesic currents}

The idea is to reconstruct a current, i.e. a cross-ratio from the periods. (talk/cite previous work by Hamenstadt, Ledrappier, Martone-Zhang?).\\
Work of Otal \cite{Otal90:SpectreMarqueNegative} shows that a current is uniquely determined by its values on closed curves, by using boxes whose endpoints are orbits of certain geodesics.
In his argument, Otal uses lifts of a dense geodesic in the surface, but controlling the asymptotics of those is complicated.
We will instead use boxes whose endpoints are orbits of a closed geodesic $x$, easier to control. The caveat here is that
now we have to be more careful since the target geodesic current may have
delta-masses at the boundaries of the boxes. We will have to specify exactly what's included in the boundary and what is not.
This is where the definition of \emph{right handed} rectangle will come into play. \\

Then we will construct closed curves whose lengths
should correspond to these rectangles. Thus, in addition to the closed geodesic
$x$, we have 4 elements of $\pi_1$, $a,b,c,d$ corresponding to the 4 sides of
these boxes ($c$ and $d$ are opposite sides, and $a,b$ are diagonals), and we will consider limits of

  $$v(a x^n b x^m) - v(c x^n d x^m).$$

One key point here is that the first curve has a crossing that the second one doesn't have
(so the difference is positive). There will be several things to check, all
up to small additive error that goes to zero as $n,m \to \infty$:

\begin{enumerate}
\item $$v(a x^n b x^m) \sim  v(a x^n) + v(b x^m)$$ (same for $c,d$).

\item Both quantities $v(a x^n b x^m)$ and $v(a x^n) + v(b x^m)$ are convex as a function of $n,m$

\item Both quantities depend only on $n+m$
\end{enumerate}


 We use then the additivity of the function $v$ to show it's additive on these boxes, and prove that the boxes form a semi-ring (see proposition \ref{prop:RHsemiring}).
 Then we apply Caratheodory's extension theorem for pre-measures on semi-rings (see theorem \ref{thm:caratheodorymeasures}) to obtain a $\pi_1(S)$-invariant measure.
 Finally, we show these approximations converge to the right value, i.e. to the value of the function on closed curves i.e. the current we have
constructed is really the Liouville current for $v$ (see proposition \ref{prop:liouvillecurrent}).

\subsection{Boxes}
\label{sec:boxes}

Let $v\colon \mathbb{R}\mathcal{C}_+(\Sigma) \to \mathbb{R}_{0 \geq}$ satisfy additive union, smoothing, stability and homogeneity.

We will fix (arbitrarily) a simple closed oriented curve~$x$ that will
serve as endpoints of the allowable rectangles. For convenience, we
will work in $\pi_1(S,\ast)$ with a basepoint~$\ast$ on~$x$, and also
use $x$ for the corresponding element of~$\pi_1$.

We say that $a \in \pi_1(S, \ast)$ \emph{crosses $x$ to the
  right} if (there is an essential right-handed crossing).

For $a,b,c,d \in \pi_1(S,\ast)$ curves that cross $x$ to
the right with $a b^{-1} c d^{-1} = 1$,
let the box $B_{a,b,c,d,x}$
be the \emph{right handed rectangle} defined by 
(show picture of lifts of curves here.)
(Need assumption on  $b$, $c$ crossing.)

Slight issue: Need a little more data to get a box in $\wt S$, rather
than a box up to translation by $\pi_1(S)$. Do this by picking
paths $p_1$, $p_2$, $q_1$, $q_2$ and setting
\begin{align*}
  a &= q_1p_1 & b &= q_2p_1\\
  c &= q_1p_2 & d &= q_2p_2.
\end{align*}

Explain about whiskers, pointing down or to the right. Allowable boxes
have all corners at whiskers.

\subsection{Measures of boxes}

We wish to define a pre-measure by setting
\todo{Labels have been permuted}
\[
\mu_{v}(B_{a,b,c,d,x})=\lim_{n \to \infty}v(bx^n) + v(cx^n)-v(ax^n) - v(dx^m)
\]
To check that this is well-defined, we first prove that this limit
exists.

\begin{lemma}
  If $a$ crosses $x$ to the right, then we have the following
  geometric reductions:
  \begin{align}
    \label{eq:cross-red1}[a x^n][x^m] &\reducesto [a x^{n+m}]\\
    \label{eq:cross-red2}[a x^{n+1}][a x^{n-1}] &\reducesto 2[a x^n]\\
    \label{eq:cross-red3}[a x^n][a x^m] & \reducesto x^{\abs{n-m}}.
  \end{align}
\end{lemma}
\begin{proof}
  Equation~(\ref{eq:cross-red1}) follows since $a$ (and $a x^n)$
  cross~$x$. The other two follow from the bands model.
\end{proof}

\begin{corollary}\label{cor:additivity-limit}
  If $a$ crosses $x$ to the right, then
  \[
    \lim_{n \to \infty} \bigl(v(a x^n) - v(x^n)\bigr)
  \]
  exists.
\end{corollary}

\begin{proof}
  By Equation~(\ref{eq:cross-red2}), $v(a x^n)$ is a convex function
  of $n$, and thus so is $v(a x^n) - v(x^n)$. By
  Equations~(\ref{eq:cross-red1}) and~(\ref{eq:cross-red3}) (both with
  $n=0$),
  \[
    -v(a) \le v(a x^n) - v(x^n) \le v(a).
  \]
  A convex function that is bounded as $n \to \infty$ has a
  well-defined limit.
\end{proof}

Thus the limit defining $\mu_{v}(B_{a,b,c,d,x})$ exists. To prove
useful facts about this limit, we need a little more.

\begin{lemma}\label{lem:join-split-geom}
  If $a$ and $b$ cross $x$ to the right, then there is an integer $k$
  so that, for all $n,m,r$ with either $0 < r < m-n+k$ or $0 > r > m-n+k$,
  \begin{align*}
    [a x^n] [b x^m] &\reducesto [ax^{n+r}bx^{m-r}]\\
    [a x^nb x^m] &\reducesto [ax^{n+r}][bx^{m-r}].
  \end{align*}
\end{lemma}
\begin{proof}
  Use the band model, centered on~$x$.
  Number the windows on top and bottom (cut out by other lifts of~$x$)
  with an arbitrary basepoint on top, and chosen so that lifts of $a$ connect
  top window~$\ell$ to bottom window~$\ell$. Choose $k$ so that lifts
  of~$b$ connect top window~$\ell$ to bottom window $\ell+k$.

  For the first reduction, note
  lifts of $a x^n$ connect windows with a shift of~$n$, and
  similarly lifts of $b x^m$ connect windows with a shift of $k+m$.
  Choose a lift of $a x^n$ that connects window $0$ to window~$n$,
  and choose a lift of $b x^m$ that connects window $n+r-k-m$ to
  window $n+r$. Smoothing gives the desired reduction.

  Other reduction is similar, instead choosing two different lifts of
  $a x^n b x^m$.
  
  See DPT notes \texttt{2018-03-01 \#1} (but it's not written up very
  cleanly there).
\end{proof}

\begin{proposition}\label{prop:join-split-alg}
  If $a$ and $b$ cross $x$ to the right, then
  \[
    \lim_{n,m\to \infty} v(a x^n b x^m) - v(ax^n) - v(bx^m) = 0.
  \]
\end{proposition}
\begin{proof}
  Let $k$ be the value from Lemma~\ref{lem:join-split-geom} for $a$ and
  $b$. Then, by that Lemma,
  \[
    v(a x^n) + v(b x^{n+k+4}) \ge v(a x^{n+1} b x^{n+k+3})
      \ge v(a x^{n+2}) + v(b x^{n+k+2}).
  \]
  In particular,
  $\lim_{n \to \infty} v(a x^n b x^{n+k+2}) - v(x^{2n+k+2})$ exists
  (and equals what it is supposed to equal).

  Muck around a little to show that the limit exists in general.
\end{proof}

\begin{proposition}
  The limit defining $\mu_{v}(B_{a,b,c,d,x})$ exists and is positive.
\end{proposition}

\begin{proof}
  Existence is clear from Corollary~\ref{cor:additivity-limit}.

  For positivity, show that $[b x^n c x^m] \reducesto [a x^m][d x^n]$.
  Apply Proposition~\ref{prop:join-split-alg}.
\end{proof}

We will address additivity of $\mu_v$ more later, but for now we prove
one easy case.

\begin{proposition}
  If two boxes $B_1$ and $B_2$ fit together side-by-side, then
  $\mu_v(B_1 \cup B_2) = \mu_v(B_1) + \mu_v(B_2)$.
\end{proposition}

\begin{proof}
  Define the boxes by elements $p_1,p_2,q_1,q_2,q_3 \in \pi_1(S, *)$.
  Then
  \begin{align*}
    B_1 &= (q_1 p_1, q_2 p_1, q_1 p_2, q_2 p_2)\\
    B_2 &= (q_2 p_1, q_3 p_1, q_2 p_2, q_3 p_2)\\
    B_1 \cup B_2 &= (q_1 p_1, q_3 p_1, q_1 p_2, q_3 p_2).
  \end{align*}
  Additivity follows immediately from the formula.
\end{proof}

\begin{remark}
Note how crucial here was that in each application of smoothing, the function satisfies the strict smoothing property. Otherwise, the asymptotics would be off by a constant.
\end{remark}

\subsection{Carathéodory's Theorem}
\label{sec:caratheodry}

\begin{definition}
semi-ring in measure theory
\end{definition}

\begin{proposition}
The right handed rectangles $B_{a,b,c,d,x}$ form a semi-ring.
\end{proposition}

(Proof is showing more: that we only need side-by-side additivity.
What's the term for that?)

\begin{proof}
Show finite additivity by using whiskers argument. Idea:
 Induction on the number of vertical+horizontal
coordinates that appear. Look at any vertical division whose whisker
reaches
the upper boundary. If there isn't one, there's a rectangle that
reaches clear across the top; remove that rectangle, by induction
build up the rest using side-by-side additivity, and then add that
missing rectangle.

Otherwise, divide the entire collection of rectangles at the location
of that vertical division. (This is possible, since this is the
correct direction for the whisker.) By induction, build up each half
by side-by-side additivity, and then attach the two sides together.
\end{proof}

\begin{definition}
pre-measure
\end{definition}

By results above, have finite additivity. Still need to prove $\mu_v$ is
countably additive.

\begin{remark}
$\mu_{\nu}$ is $\pi_1(S)$ invariant, since by definition its definition does not depend on the lift of the curves considered.
\end{remark}

\begin{theorem}
Carathéodory's theorem for semi-rings.
\end{theorem}

\subsection{Countable additivity}
\label{sec:countable-additivity}

Define continuity from above. Continuity from above is enough for
rings.

\begin{definition}
  Fix a ring $\mathcal{R}$. Let $\mathcal{F}$ be a family of nested sequences
  of elements of $\mathcal{R}$. Then $\mathcal{F}$ is \emph{final} if,
  for any sequence $(R_i)_{i=0}^{\infty}$ of nested elements of
  $\mathcal{R}$ with $\bigcap_i R_i = \emptyset$, there is a sequence
  $(S_j) \in \mathcal{F}$ so that $R_i$ is eventually contained in
  $S_j$, in the sense that for each $j$ there is an $N$ so that for
  all $i \ge N$ we have $R_i \subset S_j$.
\end{definition}

To show countable additivity, suffices to show that there is a final
family on which measures go to zero.

For $a,c,e$ crossing $x$ to the right, define a family
\[
  F_{a,b,e} = \bigl(B_{a,bx^ke,c,dx^k e}\bigr)_{k=0}^\infty
\]
\begin{itemize}
\item $F_{a,b,e}$ is nested
\item $\lim_{k \to \infty} \mu_v(F_{a,b,e}(k)) = 0$.
\item $F_{a,b,e}$ is final.
\end{itemize}

Completes proof that we have a measure.

\DMGtodo[inline]{This is the part missing to prove countable subadditivity}

\subsection{Recovering original measure}
\label{sec:recovery}

\begin{proposition}
For $C$ a closed curve, we have
\[
i(\mu_{\nu},C)=\nu(C).
\]
\end{proposition}


\subsection{The proof, older version}
\DMGtodo[inline]{This is the old version of the proof}
\begin{proposition}\label{prop:premeasurewelldef}
$\mu_{v}(B_{a,b,c,d,x})$ is well-defined.
\end{proposition}
\begin{proof}
For the elements $x,a,b,c,d \in \pi_1(S)$, we consider the following curves

\textbf{Type 1}

\begin{enumerate}
\item $cx^ndx^m$
\item $cx^n$
\item $dx^m$
\end{enumerate}

\textbf{Type 2}

\begin{enumerate}
\item $ax^nbx^m$
\item $ax^n$
\item $bx^m$
\end{enumerate}


If we define $f(n)\coloneqq v(cx^n)$, $g(m)\coloneqq v(dx^m)$ and $h(n,m)\coloneqq v(cx^ndx^m)$.
Application of smoothing and additive union yields

\begin{enumerate}
\item $f(n) \sim K_1 + nv(x)$
\item $g(m) \sim K_2 + mv(x)$.
\end{enumerate}

where $\sim$ here means asymptotically.
(show intersections are essential by working in the band model).\\
We want to show that

$h(n,m) \sim K_1 + K_2 + (n+m)v(x)$.

To do this, we use the smoothing property again, and take lifts in the bands model to see that intersections are essential and thus smoothing yields
\[
  [cx^n][dx^{n+k}] \searrow [cx^ldx^{n+k-l}]
\]
where $1 \leq l \leq k-1$.

This shows that
\[
  f(n+k) + g(n) \geq h(n+l,n+k-l).
\]
Similarly,
\[
  [cx^ndx^m][cx^k] \searrow [dx^{n+m-k}]
\]
which yields
\[
  h(n,m) + f(k) \geq g(n+m-k)
\]
i.e.,
\[
  h(n,m) \geq K_1 + K_2 + (n+m)\nu(x) - C(k).
\]
For, $n+m \gg k$ (or $n+m > k + C$). \\

Now we show that $v(cx^n)$ is convex as a function of $n$.
\end{proof}

\begin{definition}
semi-ring in measure theory
\end{definition}

\begin{proposition}\label{prop:RHsemiring}
The right handed rectangles $B_{a,b,c,d,x}$ form a semi-ring.
\end{proposition}
\begin{proof}
Show finite additivity by using whiskers argument. Idea:
 Induction on the number of vertical+horizontal
coordinates that appear. Look at any vertical division that reaches
the upper boundary. If there isn't one, there's a rectangle that
reaches clear across the top; remove that rectangle, by induction
build up the rest using side-by-side additivity, and then add that
missing rectangle.

Otherwise, divide the entire collection of rectangles at the location
of that vertical division. (This is possible, since this is the
correct direction for the whisker.) By induction, build up each half
by side-by-side additivity, and then attach the two sides together.
\end{proof}

\begin{definition}
pre-measure
\end{definition}

\begin{proposition}\label{prop:premeasure}
$\mu_{\nu}$ is a pre-measure on the semi-ring formed by the right handed rectangles.
\end{proposition}
\begin{proof}
Show *countable* additivity of pre-measure on disjoint union of right handed rectangles. Use additivity of $\nu$.
\end{proof}

\begin{remark}
$\mu_{\nu}$ is $\pi_1(S)$ invariant, since by definition its definition does not depend on the lift of the curves considered.
\end{remark}
\begin{theorem}\label{thm:caratheodorymeasures}
caratheodory's theorem for semi-rings.
\end{theorem}

\begin{proposition}
$\mu_{\nu}$ extends to a $\pi_1(S)$-invariant measure.
\end{proposition}
\begin{proof}
By \ref{prop:premeasure} and \ref{prop:RHsemiring} $\mu_{\nu}$ is a $\pi_1(S)$-invariant pre-measure defined on a semi-ring, thus by theorem \ref{thm:caratheodorymeasures} extends to a $\pi_1(S)$-invariant measure.
\end{proof}

\begin{remark}
Note how crucial here was that in each application of smoothing, the function satisfies the strict smoothing property. Otherwise, the asymptotics would be off by a constant.
\end{remark}

\begin{proposition}\label{prop:liouvillecurrent}
For $C$ a closed curve, we have
\[
i(\mu_{\nu},C)=\nu(C).
\]
\end{proposition}
\begin{proof}

\end{proof}
%%% Local Variables:
%%% mode: latex
%%% TeX-master: "Smoothings"
%%% End:

\section{Applications}

\DMGtodo[inline]{some possible applications of this result. If the conditions give us intersection numbers, then one should be able to define geodesic currents as a subset of functionals on curves. This allows to give a new definition of Thurston distance in terms of supremums over a subfamily of functionals. Then one can especulate things like: can Teichmueller distance be expressed as supremum over a larger family of functionals on curves?}

\subsection{Redefining Thurston metric}

$d_{\mathit{Th}}(X,Y)=\sup_{f \in \mathcal{C}}\frac{f(X)}{f(Y)}$ (it is not clear if $\mathcal{C}$ here should mean only additive convex functionals (motivation: true for intersection numbers).
\DMGtodo[inline]{What is the smallest family of functionals $\mathcal{C}$ that guarantees this is actually a symmetric distance? }


\subsection{Characterizing binding currents?}

$\operatorname{Syst}(\mu)=0$ if and only if $f(\mu)=0$ for some convex (additive?) functional. (motivation: result by M. Burger on structural theorem for geodesic currents).

\subsection{Anosov representations}

Let 
\[
\partial \Gamma^{[4]}\coloneqq\{(a,b,c,d) \in \partial \Gamma^{4}\colon \{ a,b \} \cap \{ c,d\}=\emptyset \}
\]

\begin{definition}
  A cross-ratio is a continuous function $B\colon \Gamma^{[4]} \to \mathbb{R}$ invariant under the diagonal action of $\pi_1(S)$ and satisfying:
  \begin{itemize}
      \item (symmetry) $B(x,y,z,w)=B(z,w,x,y)$
      \item (additivity) $B(x,y,z,w) + B(x,y,w,u)=B(x,y,z,u)$ for all $x,y,z,w,u \in \partial \Gamma$ such that $(x,y,z,w), (x,y,z,u) \in  \Gamma^{[4]}$
  \end{itemize}
\end{definition}
We say that $B$ is \emph{positive} if for all $x,y,z,w$ in this cyclic order, we have $B(x,y,z,w) \geq 0$. \\
Recall that if $c=[\gamma]$ we can define the length 
\[
  l_{\rho}^{\alpha}(c)=B_{\rho}^{\alpha}(\gamma^+,\gamma^-,z,\gamma(z)).
 \]

The properties of cross ratio ensure this is well-defined.
In \cite{MZ18:PositiveReps} it is shown that if $B_{\rho}^{\alpha}$ is positive, then there exists a unique
geodesic current $\mu_{\rho}^{\alpha}$ so that $l_{\rho}^{\alpha}(c)=i(\mu_{\rho}^{\alpha},c)$. In particular, $l_{\rho}^{\alpha}$ extends to geodesic currents as an intersection number.
We can give another characterization of positivity for cross-ratio in terms of the length function $l_{\rho}^{\alpha}$.
 
\begin{proposition}
 $B_{\rho}^{\alpha}$ is positive if and only if $l_{\rho}^{\alpha}$ satisfies smoothing condition.
\end{proposition}
\begin{proof}
From the above mentioned result in \cite{MZ18:PositiveReps}, the smoothing property follows. \\
On the other hand, if $l_{\rho}^{\alpha}$ satisfies smoothing, since additive union is also clear for $l_{\rho}^{\alpha}$, then we just apply theorem \ref{thm:intersect}
to obtain a geodesic current $\mu_{\rho}^{\alpha}$. Finally, we define the cross-ratio 
\[ 
B_{\rho}^{\alpha}(a,b,c,d)\coloneqq \mu_{\rho}^{\alpha}([a,b)\times[c,d)).
\]

It obviously satisfies the properties of cross-ratio in \cite{MZ18:PositiveReps} because of the definition of geodesic current, and it's positive.

\end{proof}
\section{Notation}


%%% Local Variables:
%%% mode: latex
%%% TeX-master: "Smoothings"
%%% End:
 % For our reference
%\nocite{*}
\bibliographystyle{hamsalpha}
\bibliography{Intersections}


\end{document}
